\documentclass[11pt]{article}
\usepackage[margin=1in]{geometry}
\usepackage{amsmath,amssymb}
\usepackage[hidelinks]{hyperref}
\emergencystretch=3em

\title{Literature-Implied Answer to a Dotted Subideal Inclusion in arXiv:1507.03241}
\author{Run fa\_banach\_001, lane 6}
\date{June 28, 2026}

\begin{document}
\maketitle

\section*{Status}

This note records a later-literature answer to the dotted-arrow question in
Section 5 of Wallis, arXiv:1507.03241. The answer comes from Freeman,
Schlumprecht, and Zsak, arXiv:1612.01153.

The packet is a literature-status record, not a new theorem. The relationship
is agent-identified from the supporting theorem and the remark following its
proof, so the status is \emph{literature-implied answer}.

\section*{Source Signal}

Wallis summarizes the then-known closed subideal structure of
$\mathcal L(\ell_p,c_0)$, $1<p<\infty$, by the diagram
\[
\{0\}\Rightarrow \mathcal K \Rightarrow
[\mathcal G_{I_{p,0}}] \ \cdots\!\!\!>\  \mathcal{FSS}
--\!\!\!>\mathcal L .
\]
The source explicitly says that a dotted arrow denotes an inclusion which is
not known to be proper.

The same section then proves a stronger diagram only for the special range
$1<p<2$, inserting further ideals between $[\mathcal G_{I_{p,0}}]$ and
$\mathcal{FSS}$. The dual statement gives the analogous information for
$\mathcal L(\ell_1,\ell_q)$ when $2<q<\infty$.

\section*{Supporting Theorem}

Freeman, Schlumprecht, and Zsak prove in arXiv:1612.01153 that for every
$1<p<\infty$ the space $\mathcal L(\ell_p,c_0)$ has continuum many closed
ideals. In Theorem 1, they construct continuum many closed ideals between
$\mathcal J^{I_{U,V}}$ and $\mathcal{FSS}$.

The remark following the proof of Theorem 1 identifies the base ideal in the
range $2\leq p<\infty$: in that range, $\mathcal J^{I_{\ell_p,c_0}}$ and
$\mathcal J^{I_{U,V}}$ are identical. Hence their continuum family lies above
$[\mathcal G_{I_{p,0}}]$ and below $\mathcal{FSS}$, proving that the dotted
inclusion from Wallis's diagram is proper for $2\leq p<\infty$.

Combining this later result with Wallis's own $1<p<2$ result answers the
dotted-inclusion question for all $1<p<\infty$.

\section*{Scope}

This note does not address the open problems left at the end of
arXiv:1612.01153, including whether the lattice embedding preserves meets or
whether $\mathcal{FSS}\subset \mathcal J^K$ in their Theorem 2.

\section*{References}

\begin{enumerate}
\item B. Wallis, \emph{Closed ideals in $\mathcal{L}(X)$ and
$\mathcal{L}(X^*)$ when $X$ contains certain copies of $\ell_p$ and $c_0$},
arXiv:1507.03241.
\item D. Freeman, Th. Schlumprecht, and A. Zsak,
\emph{Closed ideals of operators between the classical sequence spaces},
arXiv:1612.01153.
\end{enumerate}

\end{document}
